\documentclass[11pt]{article}
\usepackage[a4paper,margin=1.1in]{geometry}
\usepackage[T1]{fontenc}
\usepackage[utf8]{inputenc}
\usepackage[english]{babel}

\usepackage{lmodern}
\usepackage{amsmath,amssymb,amsthm,mathrsfs}
\usepackage{microtype}
\usepackage{fancyhdr}
\usepackage{enumitem}
\usepackage{longtable}
\usepackage{array}
\setlength{\parindent}{1.5em}
\setlength{\parskip}{0.18em}
\pagestyle{fancy}
\fancyhf{}
\cfoot{\thepage}
\newcommand{\Om}{\Omega}
\newcommand{\Omm}{\Omega_m}
\newcommand{\Hm}{H_m}
\newcommand{\Norm}{\mathrm{N}}
\newcommand{\Normone}{\mathrm{N}_1}
\newcommand{\frp}{\mathfrak p}
\newcommand{\frP}{\mathfrak P}
\newcommand{\frG}{\mathfrak G}
\newcommand{\frA}{\mathfrak A}
\newcommand{\frB}{\mathfrak B}
\newcommand{\Gg}{\mathfrak G}
\newcommand{\Aa}{\mathfrak A}
\newcommand{\Bb}{\mathfrak B}
\newcommand{\pp}{\mathfrak p}
\newcommand{\eps}{\varepsilon}
\newcommand{\srcpage}[1]{\par\medskip\noindent\textsf{[source p.~#1]}\par\smallskip}

\lhead{Heinrich Weber, Lehrbuch der Algebra, Vol. II}
\rhead{literal repair cumulative}
\begin{document}
\section*{\S\,176. The real field $H_m$ contained in $\Omega_m$.}

Every number of the field $\Omega_m$ is transformed by the substitution $(r,r^{-1})$ into the conjugate-imaginary number which is also obtained by the interchange $(i,-i)$. The product of two such conjugate-imaginary numbers is the square of the absolute value of either of them.

A number in $\Omega_m$ is real if it remains unchanged under the interchange $(r,r^{-1})$, and only under this condition. Every such number can be expressed rationally by the two-term period $r+r^{-1}$, and therefore belongs to a real field of degree $\frac12\varphi(m)$ which is a divisor of $\Omega_m$. This field we shall call the real field contained in $\Omega_m$ and denote it by $H_m$ (although other real fields are also contained in $\Omega_m$, namely all divisors of $H_m$).

The $\frac12\varphi(m)=\frac12\mu$ numbers
\begin{equation}
1,\quad r+r^{-1},\quad r^2+r^{-2},\ldots,\quad r^{\frac12\mu-1}+r^{-\frac12\mu+1}
\tag{1}
\end{equation}
form a basis of the algebraic integers of the field $H_m$. For by \S\,170, (7),
\begin{equation}
\begin{aligned}
\omega={}&x_0+x_1r+x_2r^2+\cdots+x_{\frac12\mu-1}r^{\frac12\mu-1}+x_{\frac12\mu}r^{\frac12\mu} \\
&+x_{-1}r^{-1}+x_{-2}r^{-2}+\cdots+x_{-\frac12\mu+1}r^{-\frac12\mu+1}
\end{aligned}
\tag{2}
\end{equation}
is an algebraic integer of the field $\Omega_m$ if and only if the coefficients $x_0,x_1,
\ldots$ are rational integers.

The condition that this number should belong to the field $H_m$ is obtained by carrying out the substitution $(r,r^{-1})$ and setting the difference equal to zero. If one then also divides by $r-r^{-1}$, one obtains
\begin{align*}
0={}&(x_1-x_{-1})+(x_2-x_{-2})(r+r^{-1})+\cdots \\
&+(x_{\frac12\mu-1}-x_{-\frac12\mu+1})\frac{r^{\frac12\mu-1}-r^{-\frac12\mu+1}}{r-r^{-1}}
+x_{\frac12\mu}\frac{r^{\frac12\mu}-r^{-\frac12\mu}}{r-r^{-1}}.
\end{align*}
The divisions can be carried out here; and if one then multiplies by $r^{\frac12\mu-1}$, there results for $r$ a rational equation of degree lower than the $\mu$th, which can be satisfied only if all its coefficients vanish. This leads to the equations
\[
x_{\frac12\mu}=0,\qquad x_{\frac12\mu-1}=x_{-\frac12\mu+1},\ldots,\qquad x_1=x_{-1},
\]
and hence every algebraic integer $\omega$ of the field $H_m$ is contained in the form
\[
\omega=x_0+x_1(r+r^{-1})+x_2(r^2+r^{-2})+\cdots+x_{\frac12\mu-1}(r^{\frac12\mu-1}+r^{-\frac12\mu+1}).
\]
Instead of the basis (1), if one sets
\[
\rho=r+r^{-1},
\]
one may also choose the powers of $\rho$:
\begin{equation}
1,\quad \rho,\quad \rho^2,\ldots,\quad \rho^{\frac12\mu-1}
\tag{3}
\end{equation}
as a basis of the algebraic integers of $H_m$. For the quantities (3) can be expressed integrally in terms of the quantities (1), and conversely.

The number $\rho$ is the root of an irreducible equation of degree $\frac12\mu$, which we shall denote by $\psi(\rho)=0$. If $f(r)=0$ is the equation for $r$ (\S\,169), then the identical relation
\begin{equation}
f(x)=x^{\frac12\mu}\psi\left(x+\frac1x\right)
\tag{4}
\end{equation}
holds; differentiating gives
\begin{equation}
f'(r)=r^{\frac12\mu}\psi'(\rho)(1-r^{-2}).
\tag{5}
\end{equation}
From this we can form the fundamental number $\Delta_1$ of the field $H_m$, which, since $H_m$ contains only real numbers, must be positive. If we denote by $\Normone$ the norm with respect to $H_m$, this fundamental number is
\[
\Delta_1=\pm\Normone\psi'(\rho).
\]
But if $\Norm$ is the norm formed with respect to $\Omega_m$, then
\[
\Norm\psi'(\rho)=[\Normone\psi'(\rho)]^2=\Delta_1^2,
\]
and therefore formula (5), with regard to \S\,169, gives
\[
\Delta=\Norm(1-r^{-2})\Delta_1^2.
\]
By \S\,169, (7) and (10),
\begin{equation}
\begin{aligned}
\Norm(1-r^{-2})&=q &&\text{for odd }q,\\
&=4 &&\text{for }q=2,\\
\pm\Delta&=q\Delta_1^2 &&\text{for odd }q,\\
&=4\Delta_1^2 &&\text{for }q=2.
\end{aligned}
\tag{6}
\end{equation}
If one substitutes therein
\[
\Delta=\pm q^{q^{\kappa-1}[\kappa(q-1)-1]},
\]
it follows that
\begin{equation}
\begin{aligned}
\Delta_1&=q^{\,q^{\kappa-1}\frac{q-1}{2}-\frac{q^{\kappa-1}+1}{2}} &&\text{for odd }q,\\
&=2^{(\kappa-1)2^{\kappa-2}-1} &&\text{for }q=2,
\end{aligned}
\tag{7}
\end{equation}
so that $\Delta_1$ is always a power of $q$.

\section*{\S\,177. The prime ideals in the field $H_m$.}

We now have to examine the question of the relation between the prime ideals of the field $H_m$ and those of the field $\Omega_m$.

If $\frP$ denotes any prime ideal in $H_m$ of degree $f_1$, then $\frP$ must be a divisor of a natural prime number $p$, and therefore (by \S\,171) $\frP$ can be divisible by the second or a higher power of a prime factor $\frp$ in $\Omega_m$ only if $p=q$.

If, however, $\frP$ is a divisor of $q$, then it must be a power of $\sigma=1-r$. Put, for instance,
\[
\frP=\sigma^\lambda.
\]
Taking the norm of this with respect to $\Omega_m$, which is the square of the norm with respect to $H_m$, one obtains (by \S\,169)
\[
q^{2f_1}=q^\lambda,
\]
and hence $f_1=\frac12\lambda$. Since $f_1$ is an integer, $\lambda$ must be at least equal to $2$. But $\lambda$ also cannot be greater than $2$, since $\sigma^2$ is associated with the algebraic integer $(1-r)(1-r^{-1})$ contained in $H_m$, and therefore $\sigma^2$ must be divisible by $\frP$. Consequently $f_1=1$, and we obtain the theorem:

\begin{enumerate}[label=\arabic*.,leftmargin=2.2em]
\item The prime number $q$ is the $\frac12\mu$th power of a prime number of first degree existing in $H_m$.
\end{enumerate}

Now let $p$ be different from $q$, and let $\frp$ be a prime factor of $\frP$ in the field $\Omega_m$, of degree $f$, which is carried by the substitution $(r,r^{-1})$ into $\frp'$. Then $\frP$ is divisible both by $\frp$ and by $\frp'$; on the other hand, $\frp\frp'$ is contained in $H_m$ and is therefore divisible by $\frP$. Thus we have to distinguish whether $\frp$ and $\frp'$ are distinct or not.

If $\frp$ is distinct from $\frp'$, then $\frP$ is divisible by $\frp\frp'$, and it follows that $\frP=\frp\frp'$; taking norms gives $f_1=f$.

If, however, $\frp=\frp'$, then $\frP=\frp$, and taking norms gives
\[
2f_1=f.
\]
Which of these two cases occurs depends, therefore, on whether the prime ideal $\frp$ admits the substitution $(r,r^{-1})$ or not, that is, whether $(r,r^{-1})$ is contained in the group to which the prime ideal belongs or not. By \S\,172 this difference is determined by whether there is some exponent $h$ for which the congruence
\[
p^h\equiv -1\pmod m
\]
is satisfied, or whether there is no such exponent.

We summarize the result as follows:

\begin{enumerate}[label=\arabic*.,leftmargin=2.2em,start=2]
\item The natural prime numbers $p$ different from $q$ split into two kinds, $p_1,p_2$.

The first kind $p_1$ is characterized by the condition that for some exponent $h$
\[
p_1^h\equiv -1\pmod m,
\]
and these prime numbers, if they belong to the exponent $f$ modulo $m$ and if $\mu=ef$, split in the field $H_m$ into $e$ prime factors of degree $\frac12f$. Their prime factors are not further decomposable in $\Omega_m$ either.

The prime numbers of the second kind $p_2$ are determined by the condition that none of their powers is congruent to $-1$ modulo $m$; they split in the field $H_m$ into $\frac12e$ prime factors of degree $f$. Their prime factors are still decomposable in $\Omega_m$ into two factors each.
\end{enumerate}

For prime numbers of the first kind, $f$ must certainly be even. If $m$ is odd, it is also sufficient, in order that $p$ be a prime number of the first kind, that it belong to an even exponent; for then
\[
(p^{\frac12 f}-1)(p^{\frac12 f}+1)\equiv0\pmod m.
\]
The first of these factors, $p^{\frac12 f}-1$, is not divisible by $m$, since otherwise $p$ would belong to the exponent $\frac12f$. Consequently $p^{\frac12 f}+1$ must be divisible by $q$. But here both factors $p^{\frac12 f}-1$ and $p^{\frac12 f}+1$ cannot be divisible by $q$, since otherwise their difference, which is equal to $2$, would also be divisible by $q$. Hence $p^{\frac12 f}+1$ is divisible by $m$.

If, however, $m$ is a power of $2$, then the congruence
\[
p^h\equiv -1\pmod m
\]
is possible only if $p\equiv -1\pmod m$. For every even power of an odd number is $\equiv+1\pmod 8$ and therefore cannot be $\equiv-1\pmod m$. But if $h$ is odd, then
\[
\frac{p^h+1}{p+1}=p^{h-1}-p^{h-2}+\cdots+1
\]
is itself odd, and consequently $p^h+1$ is divisible by no higher power of $2$ than $p+1$. Thus we may add to Theorem 2 the following more precise determination:

\begin{enumerate}[label=\arabic*.,leftmargin=2.2em,start=3]
\item If $m$ is odd, then the prime numbers $p_1$ belong to an even exponent, and the prime numbers $p_2$ to an odd exponent.

If $m$ is even, then the prime numbers of the form $km-1$ belong to the first kind, and all other odd prime numbers to the second kind.
\end{enumerate}

\section*{\S\,178. The units of the field $H_m$.}

The units belonging to the field $H_m$ are of special importance for the applications. The units of the field $\Omega_m$ may also be reduced to them, according to the following theorem:

\begin{enumerate}[label=\arabic*.,leftmargin=2.2em]
\item Every unit of the field $\Omega_m$ is the product of a root of unity and a real unit of the field $H_m$.
\end{enumerate}

Denote any unit of the field $\Omega_m$ by $\mathscr E(r)$. Then $\mathscr E(r^{-1})$ is the conjugate-imaginary value, which likewise represents a unit, and the quotient $\mathscr E(r):\mathscr E(r^{-1})$ is again a unit whose absolute value
\begin{equation}
\sqrt{\frac{\mathscr E(r)}{\mathscr E(r^{-1})}\frac{\mathscr E(r^{-1})}{\mathscr E(r)}}
\tag{1}
\end{equation}
is equal to $1$. Consequently this quotient is a root of unity and hence is $=\pm r^h$, where $h$ is some integral exponent (\S\,175, Theorems 1 and 2). We therefore have
\begin{equation}
\mathscr E(r)=\pm r^h\mathscr E(r^{-1}).
\tag{2}
\end{equation}

In the rest of the proof a small distinction has now to be made according as $m$ is odd or even.

First suppose that $m$ is odd. Then in (2) we may take the exponent $h$ to be even, since, if necessary, we may replace it by $h+m$. Moreover, in formula (2) only the upper sign is admissible. For if the lower sign stood there, it would follow that
\[
\mathscr E(r)\equiv\mathscr E(1)\equiv-\mathscr E(1),\qquad 2\mathscr E(1)\equiv0\,[\mathrm{mod}\,(1-r)].
\]
But by \S\,169, $1-r$ is a divisor of $q$, and so is relatively prime to $2$; therefore $\mathscr E(1)$, and consequently also $\mathscr E(r)$, is divisible by $(1-r)$. This contradicts the supposition that $\mathscr E(r)$ is to be a unit.

Consequently (2) gives
\[
r^{-\frac12 h}\mathscr E(r)=r^{\frac12 h}\mathscr E(r^{-1}),
\]
whence it follows that $r^{-\frac12 h}\mathscr E(r)$ is a real unit. Denoting it by $e(r)$, we therefore have
\begin{equation}
\mathscr E(r)=r^{\frac12 h}e(r),
\tag{3}
\end{equation}
which proves Theorem 1 in this case.

If, secondly, $m$ is a power of $2$, then $\mu=\frac12m$ and $r^\mu=-1$; and along with $r$, $-r$ is also an $m$th root of unity. If necessary replacing $h$ by $h+\frac12m$, we may assume the upper sign in (2) and write
\begin{equation}
\mathscr E(r)=r^h\mathscr E(r^{-1}).
\tag{4}
\end{equation}
The point is now to prove that $h$ must be even.

If we represent $\mathscr E(r)$ in the basis $1,r,r^2,\ldots,r^{\mu-1}$, then
\begin{equation}
\mathscr E(r)=\sum_{s=0}^{\mu-1}x_sr^s,
\tag{5}
\end{equation}
where the $x_s$ are rational integers. But by (4) we then have
\[
\sum_{s=0}^{\mu-1}x_sr^s=\sum_{s=1}^{\mu-1}x_sr^{h-s},
\]
and from this it follows that $x_s=\pm x_{s'}$ whenever $s+s'\equiv h\pmod\mu$. If $h$ is odd, then of two numbers so connected one is even and the other odd (since $\mu$ is a power of $2$). Thus the expression (5) gives for $\mathscr E(r)$ an aggregate of terms of the form
\[
x_s(r^s\pm r^{s'}).
\]
But
\[
r^s\pm r^{s'}=r^s(1\pm r^{s'-s})=r^s(1\mp r^{\mu+s'-s})
\]
is always divisible by $1-r$, and from this it would follow that $\mathscr E(r)$ is divisible by $1-r$, which is not a unit, contrary to the concept of a unit. Therefore the assumption that the exponent $h$ in formula (4) is odd is inadmissible, and, as above, we may put
\[
\mathscr E(r)=r^{\frac12h}e(r),
\]
where $e(r)$ is a real unit. Theorem 1 is thereby proved.

We shall single out a system of real units. By \S\,169 the $\mu$ quantities $1-r^n$ are all associated with one another. Therefore the quotient
\[
\frac{1-r^n}{1-r^{n'}},
\]
where $n,n'$ are two numbers relatively prime to $m$, is a unit. From this, however, setting
\[
r=e^{\frac{2\pi i}{m}},
\]
one derives the real unit
\begin{equation}
\frac{\sin\frac{n\pi}{m}}{\sin\frac{n'\pi}{m}}
=\frac{r^{\frac n2}-r^{-\frac n2}}{r^{\frac{n'}2}-r^{-\frac{n'}2}}
=r^{\frac{n'-n}{2}}\frac{1-r^n}{1-r^{n'}}.
\tag{6}
\end{equation}
If $m$ is even, one may set $n'=n+\frac m2$, and obtains in this case the real units
\begin{equation}
\tau_n=\operatorname{tang}\frac{n\pi}{m},
\tag{7}
\end{equation}
where $n$ may denote any odd number.

Replacing $n$ by $\frac12m-n$ carries $\tau_n$ into its reciprocal value. If $n$ runs through the sequence of numbers
\[
1,3,5,\ldots,\frac m4-1,
\]
then all the values $\tau_n$ obtained are positive proper fractions.

\clearpage

\srcpage{648}
\section*{Twenty-second Section. Abelian fields and cyclotomic fields.}

\section*{\S~179. Simple Abelian fields.}

We now turn to one of the most interesting applications of the theory of algebraic numbers.

The matter is the proof of the general theorem:\footnote{Kronecker first stated this theorem, but did not publish a complete proof of it. The first proof was made known by the author of the present work in vol.~8 of \textit{Acta Mathematica} (1886). Very recently a second proof by Hilbert has been published, which is based on the theorems developed in the nineteenth section (\textit{Nachrichten d. Gesellschaft d. Wissenschaften zu Göttingen}, 1896, issue 1).}

\textbf{I. All Abelian number fields in the absolute domain of rationality are cyclotomic fields.}

Once we have proved this theorem, the investigations of the third section on cyclotomic fields acquire an increased interest, because it is thereby shown that on the path indicated there one finds not only all cyclotomic fields, but all Abelian fields whatever.

Let
\[
  \Om=R(x)
\]
be an Abelian field of degree $m$, that is, a field consisting of all rational functions of a root $x$ of an irreducible Abelian equation of degree $m$, when the field $R$ of rational numbers is regarded as the domain of rationality. The Galois group $\Gg$ of this equation, which is therefore an Abelian group and has the same degree $m$ as the field $\Om$, is also the group of the field $\Om$.

If $x_1$ is a non-primitive number from $\Om$, then $\Om_1=R(x_1)$ is likewise an Abelian field, which we call a proper divisor of $\Om$, and whose degree is a proper divisor of the degree of $\Om$. For if $\Aa$ is the group to which the number $x_1$ belongs, then $\Gg\mid\Aa$ (\S~4 of this volume) is the group of the field $\Om_1$, and this is, together with $\Gg$, an Abelian group. A divisor of $\Om$ which has the same degree as $\Om$ is identical with $\Om$ (cf. Vol.~I, \S\S~144, 149, 156).

\srcpage{649}
If there are now two non-primitive numbers $x_1,x_2$ in $\Om$ of such a kind that
\[
  \Om=R(x_1,x_2)
\]
can be put, then the fields $\Om_1=R(x_1)$, $\Om_2=R(x_2)$ are proper divisors of $\Om$, and $\Om$ is called \emph{composed} from the two fields $\Om_1,\Om_2$.

If the field $\Om$ permits no such representation, it is called \emph{simple}.

Thus, if it is assumed as proved that $\Om_1,\Om_2$ are cyclotomic fields, that is, that all their numbers are representable rationally by roots of unity, then the same follows for $\Om$.

If one of the fields $\Om_1,\Om_2$ is decomposable, then we can repeat the decomposition, and, since the degrees are decreasing positive integers, we must finally arrive at simple fields.

Theorem I therefore needs to be proved only for simple Abelian fields.

It is, however, to be recalled that the decomposability of a field $\Om$ by no means already follows from the existence of a divisor, and that in a decomposable field the decomposition into simple fields can occur in quite different ways, so that for fields the laws of divisibility do not hold as they do in the domain of numbers.

A criterion for the simplicity of a field can be derived from its group.

If the group $\Gg$ has two proper divisors $\Aa$ and $\Bb$ of such a kind that every element can be composed in one and only one way from an element of $\Aa$ and an element of $\Bb$, then we call the group $\Gg$ decomposable into the two components $\Aa,\Bb$; thus $\Gg$ is decomposable when, in the product formed according to the composition of parts (\S~4),
\begin{equation}
  \Gg=\Aa\Bb,
\tag{1}
\end{equation}
each element of $\Gg$ appears once and only once. The degree of $\Gg$ is then equal to the product of the degrees of $\Aa$ and $\Bb$.

If there are no such divisors, the group $\Gg$ is called indecomposable.

\srcpage{650}
Then we can prove the theorem:

\textbf{1. If the group of an Abelian field is decomposable, then the field too is composed.}

To prove this, assume that the group $\Gg$ of the field $\Om$ is decomposable according to formula (1), and let $m,a,b$ be the degrees of $\Gg,\Aa,\Bb$. We take a number $\xi$ of the field $\Om$ belonging to the group $\Bb$, which is therefore transformed by the substitutions of $\Gg$ into $a$ different conjugate values, and likewise a $b$-valued number $\eta$ belonging to $\Aa$. We can then form a number
\[
  x=\alpha\xi+\beta\eta
\]
with rational numerical coefficients $\alpha,\beta$, which has $m$ different conjugate values, and then $\Om=R(x)$. Thus $\Om$ is composed from $R(\xi)$ and $R(\eta)$ (Vol.~I, \S~143).

If we now recall the representation of an Abelian group by a basis derived in \S\S~9, 10 of this volume, then repeated application of what has just been proved gives:

\textbf{2. Every Abelian field $\Om$ can be composed from such fields whose degree is a power of a prime number and whose group is cyclic. The degrees of these composing fields are the invariants of the group of $\Om$.}

These theorems are supplemented by the following:

\textbf{3. An Abelian field of prime-power degree with cyclic group is simple.}

If the group $\Gg$ is cyclic, then its elements can be represented by the powers of a basis, $G^\gamma$, when $\gamma$ runs through a complete system of residues modulo $m$. We obtain all divisors $\Aa$ of $\Gg$ represented by
\[
  \Aa=G^{b\gamma},
\]
when $b$ is a divisor of $m$ and $\gamma$ runs through a complete system of residues modulo $a=m:b$.

If
\[
  \Aa'=G^{b'\gamma}
\]
is a second divisor of $\Gg$ of degree $a'$ and if
\[
  b'\le b,
  \qquad a'\ge a,
\]
then, when we now assume that $m$ and hence also $a,b,a',b'$ are powers of a prime number $q$, $b'$ is a divisor of $b$, and consequently $\Aa$ is a divisor of $\Aa'$.

\srcpage{651}
If therefore $\xi$ is a number of the field $\Om$ which belongs to the group $\Aa$, then $\xi$ is a primitive number of the field $R(\xi)$ of degree $b$, and in this field every number $\xi'$ is also contained which admits the substitutions of the group $\Aa'$. Thus if $\Aa,\Bb$ and $\Aa',\Bb'$ are proper divisors of $\Gg$, then $R(\xi)$ is different from $\Om$, and $R(\xi)$ is not enlarged by the adjunction of $\xi'$. Hence $\Om$ also cannot be composed from $R(\xi)$ and $R(\xi')$, whereby 3 is proved.

\section*{\S~180. The resolvents.}

According to what has now been proved, in the further considerations, whose aim is the proof of Theorem I, we may restrict ourselves to such Abelian fields whose degree is a prime power and whose group is cyclic. For this reason it also sufficed, for the intended application, that in the preceding section we restricted ourselves to the consideration of the cyclotomic fields $\Omega_m$ in which $m$ was a prime power. We retain as much as possible of the notation used there and put
\begin{equation}
  m=q^\kappa,
\tag{1}
\end{equation}
where $q$ is a prime number and $\kappa$ a positive exponent; and in the case where $q=2$ we assume $\kappa>1$. By $n$ we denote any number not divisible by $q$. Let $r$ be a primitive $m$-th root of unity, and $\Omega_m$ the field of rational functions of $r$.

Let now $\Om=R(x)$ be a simple Abelian field of degree $m$, so that $x$ is a root of a rational irreducible cyclic equation of degree $m$. We denote the roots of this equation, in a suitable order, by
\begin{equation}
  x_0,x_1,x_2,\ldots,x_{m-1},
\tag{2}
\end{equation}
and assume $x_m=x_0$, $x_{m+1}=x_1$, and so on. Then
\[
  x_1=\Phi(x_0),\quad x_2=\Phi(x_1),\ldots,\quad x_0=\Phi(x_{m-1}),
\]
where $\Phi$ denotes an integral function with rational coefficients, and the group of the field $\Om$ consists of the powers of the substitution $(x_0,x_1)$, or of the cyclic permutations

\srcpage{652}
of the roots (2). Every cyclic function of these quantities is a rational number. What we have to prove is that the $x$ can be expressed rationally by roots of unity. If we adjoin to the field $\Om$ the $m$-th root of unity $r$, there arises a field $R(x,r)$ which contains the two fields $\Om=R(x)$ and $R(r)=\Omega_m$ as divisors.

If a number $\omega=\Phi(x,r)$ of the field $R(x,r)$ admits all substitutions $(x_0,x_1)$, $(x_0,x_2)$, \ldots, $(x_0,x_{m-1})$, then it is contained in the field $\Omega_m$; for then
\[
  m\omega=\Phi(x_0,r)+\Phi(x_1,r)+\cdots+\Phi(x_{m-1},r),
\]
and the coefficients of the powers of $r$ in this expression are not only cyclic but even symmetric functions of the roots (2), hence rational numbers. This remains true even in the case where the equation for $x$ is reduced by adjunction of $r$.

We apply this to the resolvents
\begin{equation}
  (r,x_0)=x_0+rx_1+r^2x_2+\cdots+r^{m-1}x_{m-1},
\tag{3}
\end{equation}
and, more generally, when $s$ denotes an arbitrary integral exponent,
\begin{equation}
  (r^s,x_0)=x_0+r^s x_1+r^{2s}x_2+\cdots+r^{(m-1)s}x_{m-1}.
\tag{4}
\end{equation}
By these resolvents $x_0$ itself can again be rationally expressed, namely
\begin{equation}
  m x_0=\sum_{s=0}^{m-1}(r^s,x_0),
\tag{5}
\end{equation}
and if therefore we can show that all these resolvents $(r^s,x_0)$ are cyclotomic numbers, then we have reached our goal. If we put
\begin{equation}
  (r^s,x_\lambda)=x_\lambda+r^s x_{\lambda+1}+r^{2s}x_{\lambda+2}+\cdots+r^{(m-1)s}x_{\lambda+m-1},
\tag{6}
\end{equation}
then $(r^s,x_\lambda)$ arises from $(r^s,x_0)$ by the substitution $(x_0,x_\lambda)$, and the fundamental relation holds:
\begin{equation}
  r^{\lambda s}(r^s,x_\lambda)=(r^s,x_0).
\tag{7}
\end{equation}
All these resolvents are numbers of the field $R(x,r)$; we consider preferably the following two combinations:
\begin{align}
  (r^s,x_0)^m &= \omega_s, \tag{8}\\
  (r^s,x_0)^{-1}(r,x_0)^s &= \alpha. \tag{9}
\end{align}

\srcpage{653}
These two numbers admit, as follows immediately from (7), the substitutions $(x_0,x_\lambda)$ and are therefore numbers of the field $\Omega_m$.

If $n$ is not divisible by $q$, then, because of the irreducibility of the cyclotomic equation, none of the resolvents $(r^n,x_0)$ can vanish unless they all vanish; for the equation
\[
  (r,x_0)^m=0
\]
would be an equation for $r$ with rational coefficients, admitting all substitutions $(r,r^n)$.

Likewise, none of the resolvents $(r^s,x_0)$ can vanish unless the whole family vanishes in which $s$ has a fixed greatest common divisor with $m$.

If all $(r,x_0)$ vanish, formula (9) is not applicable. Then, however, we can directly consider the resolvents $(r^q,x_0)$, which belong to a field of degree $m q^{-1}$ arising from
\[
  y_0=x_0+x_q+x_{2q}+\cdots+x_{m-q},
\]
and if these quantities too vanish, we pass to the resolvents $(r^{q^2},x_0)$, and so on.

If, however, $(r,x_0)$ does not vanish, then by formula (9) all $(r^s,x_0)$ can be represented rationally by $(r,x_0)$ and by cyclotomic numbers.

Consequently it is sufficient for our purpose if we can prove that the non-vanishing resolvents $(r,x_0)$ are cyclotomic numbers.

Let $n$ run through a complete system of incongruent numbers not divisible by $q$. Then
\[
  \sum n\equiv 0\pmod m,
\]
for the numbers $n$ split into pairs which complement one another to $m$. From (7) it follows that the product
\begin{equation}
  \prod_n (r^n,x_0)=a
\tag{10}
\end{equation}
is a number in $\Omega_m$. But since this number moreover is not changed when any of the substitutions $(r,r^n)$ is performed in it, $a$ is a rational number.

\srcpage{654}
\section*{\S~181. Preparation for the proof.}

After the developments of the preceding paragraph, for the proof of the great Theorem I it remains only to show that the number contained in $\Omega_m$,
\[
  \omega=(r,x_0)^m,
\]
is the $m$-th power of a cyclotomic number.

\textbf{1. The essential property of the numbers $\omega_n$ conjugate to $\omega$, on which the proof rests, is that}
\begin{equation}
  \omega_n^{-1}\omega_1^n=\alpha^m
\tag{1}
\end{equation}
\textbf{is the $m$-th power of a number in $\Omega_m$.}

Here $n$ may be any number not divisible by $q$. If we form, of the right- and left-hand sides of (1), the norms in the field $\Omega_m$, then, denoting by $a$ the norm of the number $\alpha$, which is a rational number, and observing that every $\omega_n$ has the same norm as $\omega$, we obtain
\[
  N(\omega)^{n-1}=a^m.
\]
If now $m$ is odd, then we can put $n=2$ in this, and obtain as a consequence of 1:

\textbf{2. The norm of $\omega$ is the $m$-th power of a rational number,}
\begin{equation}
  N(\omega)=a^m.
\tag{2}
\end{equation}
If $m$ is a power of $2$, this inference is no longer permissible. It then follows from 1 only that $N(\omega)^2$ is the $m$-th power of a rational number. Nevertheless, as is seen from (10) of the preceding paragraph, the numbers $\omega$ must, also in this case, still satisfy condition 2; but this condition is then no longer contained automatically in 1, and is added as a new property.

What we have to prove is that from 1 and 2 it follows that $\omega$ itself is the $m$-th power of a cyclotomic number, even if in a higher field.

For this it is necessary to investigate the decomposition of the numbers $\omega$ into their prime factors in the field $\Omega_m$.

The prime number $q$ is, as was proved in \S~169, associated with the $\varphi(m)$-th power of a prime number $\sigma=1-r$ existing in $\Omega_m$.

\srcpage{655}
If therefore $\sigma^k$ is the power of $\sigma$ contained in $\omega$, we put
\[
  \omega=\sigma^k\omega',
\]
where $\omega'$ is relatively prime to the prime number $q$. Formation of the norm gives, by 2,
\[
  a^m=q^k b,
\]
where $b$ is a rational number which, in lowest terms, contains the prime number $q$ neither in the numerator nor in the denominator. From this it follows that $k$ must be divisible by $m$, and hence the first result:

\textbf{3. The exponent of the prime number $\sigma$ contained in $\omega$ is always divisible by $m$.}

Let now $p$ be a prime number different from $q$, and let $\pp_\xi$ be a prime factor of $p$ in the field $\Omega_m$. Let $\xi$ run through the sequence of numbers $\xi_1,\xi_2,\ldots,\xi_e$ determined more precisely in \S~172 in the different cases, so that
\begin{equation}
  p=\pp_{\xi_1}\pp_{\xi_2}\cdots\pp_{\xi_e},
\tag{3}
\end{equation}
and assume that $\pp_\xi^{k_\xi}$ is the power of $\pp_\xi$ contained in $\omega$. The power of $\pp_{n\xi}$ contained in $\omega_n$ is then $\pp_{n\xi}^{k_\xi}$; and if we determine $n'$ from the congruence
\begin{equation}
  nn'\equiv 1\pmod m,
\tag{4}
\end{equation}
then $\pp_\xi^{k_{n'\xi}}$ is the power of $\pp_\xi$ contained in $\omega_n$. Accordingly, in
\begin{equation}
  \omega_n^{-1}\omega_1^n
\tag{5}
\end{equation}
the power $\pp_\xi^{n k_\xi-k_{n'\xi}}$ is contained, and by 1 the exponent of this power must be divisible by $m$.

Thus for every $n$ not divisible by $q$ we obtain the congruence
\begin{equation}
  n k_\xi\equiv k_{n'\xi}\pmod m,
\tag{6}
\end{equation}
in which we can now also replace $\xi$ by $n\xi$, whereby
\begin{equation}
  n k_{n\xi}\equiv k_\xi\pmod m
\tag{7}
\end{equation}
is obtained. If in this we take $\xi=1$ and put $\xi,\xi'$ for $n,n'$ and $k$ for $k_1$, then from (7) follows
\begin{equation}
  k_\xi\equiv \xi' k\pmod m,
\tag{8}
\end{equation}
where $k$ now has the same value for all $\xi$, and $\xi,\xi'$ are connected by the congruence
\[
  \xi\xi'\equiv1\pmod m.
\]

\srcpage{656}
If we take $n=p$, then $\pp_{p\xi}=\pp_\xi$ (\S~172), and so also $k_{p\xi}=k_\xi$; and from (7) and (8) follows
\begin{equation}
  k(p-1)\equiv0\pmod m.
\tag{9}
\end{equation}
Important consequences now follow from this. First, if $p-1$ is not divisible by $q$, then it follows that $k$ must be divisible by $m$, and we therefore have the theorem:

\textbf{4. The prime factors of a prime number $p$ which is not congruent to $1$ modulo $q$ are contained in $\omega$ only in such powers whose exponents are divisible by $m$.}

If, in the case where $m$ is a power of $2$, $p-1$ is divisible by $2$ but not by $4$, that is, if $p\equiv -1\pmod4$, then it follows from (9) only that $k$ is divisible by $\tfrac12m$, not that $k$ is divisible by $m$. The exponent of the power of $\pp_\xi$ contained in $\omega$ is, apart from multiples of $m$, equal to $k$; and if $k$ is divisible by $m$, everything is as in case 4. But if $k$ is divisible only by $\tfrac12m$, not by $m$, so that
\[
  k\equiv \tfrac12m\pmod m,
\]
then $k-\tfrac12m$ is divisible by $m$, and with reference to the decomposition (3) we can formulate the theorem as follows:

\textbf{5. If $m$ is a power of $2$ and $p$ is a prime number congruent to $-1$ modulo $4$, then the exponent $h$ can be so determined that all prime factors of $p$ in}
\[
  \omega\sqrt{p}^{\,-hm}
\]
\textbf{are contained only in such powers whose exponent is divisible by $m$.}

Finally let $p-1$ be divisible by $q$, or, if $m$ is even, by $4$. Let $m_1$ be the greatest common divisor of $m$ and $p-1$, and let
\[
  m=m_1m_2.
\]
Then by (9) $k$ must be divisible by $m_2$; and if we put
\[
  k=hm_2,
\]
then by (8) the product contained in $\omega$ is
\begin{equation}
  \left(\prod_\xi \pp_\xi^{\xi'}\right)^{hm_2}.
\tag{10}
\end{equation}
In addition, the prime factors $\pp_\xi$ occur in $\omega$ only in such powers whose exponents are divisible by $m$.

\srcpage{657}
In this case $\xi$, by \S~172, III, runs through a complete system of residues not divisible by $q$ modulo $m_1$, and in (10) we may understand by $\xi'$ the least positive solution of the congruence
\[
  \xi\xi'\equiv1\pmod{m_1},
\]
because, if we change the exponents $\xi'$ in (10) modulo $m_1$, only $m$-th powers of the prime factors $\pp$ are added.

Then we can apply Kummer's theorem [\S~174, (16)] to the product (10), replacing $m$ there by $m_1$, and obtain, if we understand by $r_1=r^{m_2}$ an $m_1$-th root of unity and by the $\eta$ certain periods of the $p$-th roots of unity,
\[
  \prod_\xi \pp_\xi^{\xi'}=(r_1,\eta)^{m_1},
\]
and consequently
\begin{equation}
  \left(\prod_\xi \pp_\xi^{\xi'}\right)^{hm_2}=(r_1,\eta)^{hm}.
\tag{11}
\end{equation}
Consequently we have the theorem:

\textbf{6. If $p$ is a prime number and $p-1$ is divisible by $q$, or, when $q=2$, by $4$, then the positive exponent $h$ can be so determined that the number contained in $\Omega_m$}
\[
  \omega(r_1,\eta)^{-hm}
\]
\textbf{contains the prime factors of $p$ only in such powers whose exponent is divisible by $m$.}

It is now also no longer necessary to distinguish Theorem 5 from Theorem 6; for for $m_1=2$, 6 passes into 5, because then $r_1=-1$, and by Vol.~I, \S~171,
\[
  (r_1,\eta)^m=(-1,\eta)^m=\sqrt{p^{\,m}}
\]
holds.\footnote{In the author's paper in vol.~8 of \textit{Acta Mathematica} it was omitted to emphasize separately the case of Theorem 5.}

If we collect the result of Theorems 3 to 6 in one formula, we obtain, when $\eps$ denotes a unit functional, $\varphi$ any functional of the field $\Omega_m$, and $p$ runs through a sequence of prime numbers congruent to $1$ modulo $q$, among which the same prime number $p$ may occur several times,
\begin{equation}
  \omega=(r,x_0)^m=\eps\,\varphi^m\prod_p (r^{m_2},\eta)^m.
\tag{12}
\end{equation}

\srcpage{658}
The resolvents of the cyclotomy occurring in these formulae, $(r^{m_2},\eta)$, are cyclotomic numbers in a higher field. But their $m$-th powers are numbers contained in the field $\Omega_m$ which are transformed by the substitution $(r,r^n)$ into
\[
  (r^{n m_2},\eta)^m.
\]
Likewise the combinations
\begin{equation}
  (r^{n m_2},\eta)^{-1}(r^{m_2},\eta)^n
\tag{13}
\end{equation}
are contained in the field $\Omega_m$ [\S~174, (9), (11)]. These resolvents are special cases of the general resolvent $(r,x_0)$. From (12), if we denote by $\varphi_n,\eps_n$ the functionals conjugate to $\varphi$ and $\eps$, we obtain
\begin{equation}
  \omega_n=\eps_n\varphi_n^m\prod_p (r^{m_2n},\eta)^m.
\tag{14}
\end{equation}
This formula teaches us the most important properties of the functionals $\varphi_n$ and $\eps_n$, first:

\textbf{7. The functionals $\varphi_n^m$ are associated with numbers of the field $\Omega_m$, and therefore belong to the principal class (\S~152).}

By 1, $\omega_n^{-1}\omega_1^n=\alpha^m$ is the $m$-th power of a number of $\Omega_m$. Since the combinations (13) are also numbers in $\Omega_m$, we can put
\[
  \prod_p (r^{n m_2},\eta)^{-1}(r^{m_2},\eta)^n=\frac{\alpha}{\beta},
\]
and obtain from (14)
\[
  \eps_n^{-1}\eps_1^n(\varphi_n^{-1}\varphi_1^n)^m=\beta^m,
\]
where $\beta$ is a number in $\Omega_m$.

It follows from this that the product
\[
  \beta\varphi_n\varphi_1^{-n}=\eps
\]
is a unit functional of the field $\Omega_m$, and from this arise, for the functionals $\varphi_n$ and the unit functionals $\eps_n$, the following two theorems:

\textbf{8. The conjugate functionals $\varphi_n$ have the property that all products}
\[
  \varphi_n\varphi_1^{-n}
\]
\textbf{belong to the principal class.}

\textbf{9. The unit functionals $\eps_n$ have the property that the products}
\[
  \eps_n\eps_1^{-n}=\eps^m
\]
\textbf{are $m$-th powers of unit functionals.}


\srcpage{659}
For the proof of the principal theorem I, everything now depends on the proof of the two propositions which are to be inferred from the properties of the functionals $\varphi$ and $\eps$:

\textbf{A. The functionals $\varphi_n$ belong to the principal class.}

If we may presuppose this lemma, then $\varphi_n$ is associated with a number $\alpha_n$, and we can write formula (14), with a somewhat changed meaning of $\eps_n$, in the form
\begin{equation}
  \omega_n=\eps_n\alpha_n^m\prod_p(r^{m_2n},\eta)^m.
\tag{15}
\end{equation}
In this formula, however, $\eps_n$ is a numerical unit of the field $\Omega_m$ satisfying theorem 9.

If, therefore, from theorem 9 we can still prove the second lemma:

\textbf{B. A numerical unit $\eps_n$ of the field $\Omega_m$, which satisfies theorem 9, is the $m$-th power of a unit in $\Omega_m$, multiplied by a power of $r$;}

then in (15) we can extract the $m$-th root, and obtain, when by $\rho$ a root of unity of possibly higher degree than $m$ and by $\alpha$ a number of the field $\Omega_m$ is denoted,
\begin{equation}
  (r,x_0)=\rho\alpha\prod_p(r^{m_2},\eta),
\tag{16}
\end{equation}
where now on the right there stand nothing but cyclotomic numbers, and by this Theorem I is therefore proved.

\srcpage{659}
\section*{\S\,182. Proof of the First Auxiliary Lemma for an Odd $m$.}

In the preceding paragraph we made the proof of Theorem I dependent on two auxiliary theorems, whose proof is now still to be sought.

We formulate the first of these auxiliary theorems as follows:

\begin{enumerate}[label=\alph*),leftmargin=2.2em]
\item If $\varphi_n$ is a system of non-zero conjugate functionals of the field $\Omega_m$, of which it is known that
\[
\varphi_n^m\quad\text{and}\quad\varphi_n^{-1}\varphi_1^n
\]
belong to the principal class, then the functionals $\varphi$ themselves belong to the principal class.
\end{enumerate}

The theorem will be proved if, for some power $\varphi^k$ of $\varphi$ whose exponent is not divisible by $q$, it can be proved that it belongs to the principal class. For if $\alpha,\beta$ are numbers in $\Omega_m$ and
\[
\varphi^m=\varepsilon'\alpha,
\qquad
\varphi^k=\varepsilon''\beta,
\]
then we can determine the rational integers $x,y$ so that
\[
mx+ky=1,
\]
and then
\[
\varphi=\varphi^{mx}\varphi^{ky}=\varepsilon^{\prime\prime\prime}\alpha^x\beta^y;
\]
thus, if $\varepsilon',\varepsilon'',\varepsilon^{\prime\prime\prime}$ are units, $\varphi$ itself is associated with a number.

According to the present state of the matter, the proof for an even $m$ is to be carried out by a quite different path from that for an odd $m$; we therefore occupy ourselves first with the easier case of an odd $m$\footnote{Compare, moreover, the new proof by Hilbert, in which the difference between even and odd $m$ recedes much more (p. 642).}.

The assumptions made in a) about $\varphi_n$ can, if $\alpha,\beta$ again denote arbitrary numbers of the field $\Omega_m$ and $\varepsilon$ unit functionals, be expressed by the two formulas
\begin{align}
\varphi_n^m&=\varepsilon\beta, \tag{1}\\
\varphi_n&=\varepsilon\alpha\varphi_1^n. \tag{2}
\end{align}
We shall now solve the problem step by step.

First let, as before, $\sigma$ be the prime factor contained in $q$, which is a number existing in $\Omega_m$, and let $\sigma^k$ be the power of $\sigma$ contained in $\varphi$; put
\begin{equation}
\varphi=\sigma^k\psi,
\tag{3}
\end{equation}
where $\psi$ denotes a functional prime to $q$, which satisfies the same conditions (1), (2) as $\varphi$, and if one of the two belongs to the principal class, the same is true of the other. Thus theorem a) needs to be proved only under the assumption that $\varphi$ is prime to $q$.

Now let $p$ be a prime number different from $q$, whose decomposition into prime factors has the expression
\begin{equation}
p=\prod_\xi \frp_\xi,
\tag{4}
\end{equation}
where $\xi$ runs through the number sequence more precisely determined in \S\,172. Let $\frp_\xi^{k_\xi}$ be the power of $\frp_\xi$ contained in $\varphi$, so that, if we set
\[
\varphi=\chi\prod_\xi \frp_\xi^{k_\xi},
\]
then $\chi$ is prime to $p$ and is related to no other prime numbers (see the end of \S\,173) than $\varphi$ itself. From this there results
\begin{equation}
\varphi_n=\chi_n\prod_\xi \frp_{n\xi}^{k_\xi},
\tag{5}
\end{equation}
and in this we let $n$ again run through the sequence of the numbers $\xi$. If we further set
\[
\psi=\prod_\xi \chi_\xi,
\]
then $\psi$ too is prime to $p$ and related to no other prime numbers than $\varphi$, and from (5), using (4), we obtain
\begin{equation}
\prod_\xi\varphi_\xi=\psi p^{\sum k_\xi},
\tag{6}
\end{equation}
and from this it follows that the functional $\psi$ satisfies the same two conditions (1), (2) as $\varphi$.

It is further shown that, if the $\varphi_\xi$ lie in the principal class, then $\psi$ is also contained in the principal class.

But if $p-1$ is not divisible by $q$, then we can also conclude the converse. For by (2), $\varphi_\xi$ is equivalent to $\varphi^\xi$, and consequently
\[
\prod_\xi\varphi_\xi\ \text{is equivalent to}\ \varphi^{\sum\xi},
\qquad
\text{is equivalent to}\ \psi,
\]
and if $\psi$ is associated with a number, the same holds for $\varphi^{\sum\xi}$. By the theorem \S\,172, however, if $p-1$ is not divisible by $q$, $\sum\xi$ is also not divisible by $q$, and consequently $\varphi$ itself is contained in the principal class.

This consideration can now be repeated if $\psi$ is still related to further prime numbers which are not congruent to 1 modulo $q$, and we arrive at the theorem:
\begin{center}
Theorem a) is proved in general if it can be proved under the assumption that $\varphi$ is related only to such prime numbers as are congruent to 1 modulo $q$.
\end{center}

If we make this assumption about the functionals $\varphi_n$, then we can apply Kummer's theorem in the form of \S\,174, 3.

In what follows we denote by $\varepsilon$ unit functionals, whose more precise knowledge is immaterial, and set
\begin{equation}
\Phi_n=\prod_t \varphi_{nt'}^{t}=\varepsilon\vartheta_n^m,
\tag{7}
\end{equation}
where $t$ runs through the residues of $m$ not divisible by $q$ which are smaller than $m$, and $t'$ is determined by $tt'\equiv1\pmod m$. Then $\vartheta_n$ is a cyclotomic number in a higher field, whose $m$th power is contained in $\Omega_m$, and which satisfies the second condition, namely that
\begin{equation}
\vartheta_n^{-1}\vartheta_1^n=\alpha
\tag{8}
\end{equation}
is a number of the field $\Omega_m$.

We now introduce a frequently used sign, understanding by $E(x)$ the greatest integer contained in the real number $x$, and by $(x)$ the remainder, which is always a proper fraction and, if $x$ itself should be an integer, is to be put equal to zero. Then
\begin{equation}
x=E(x)+(x).
\tag{9}
\end{equation}
If $x$ is an integer, then in this notation
\[
m\left(\frac xm\right)
\]
is also an integer, namely the least positive remainder of the division of $x$ by $m$.

In formula (7), the index of the functionals $\varphi$ is determined only modulo $m$. The exponent, however, is reduced to its least remainder modulo $m$. Thus if in (7) we replace the index $t'$ by $n't'$, then $t$ must be replaced by the remainder of the division of $nt$ by $m$, that is, [by (9)] by
\[
m\left(\frac{nt}{m}\right)=nt-mE\left(\frac{nt}{m}\right),
\]
and we can set
\[
\Phi_n=\prod_t \varphi_{t'}^{nt-mE(\frac{nt}{m})}=\varepsilon\vartheta_n^m,
\]
whereas
\[
\Phi_1=\prod_t \varphi_{t'}^t=\varepsilon\vartheta_1^m.
\]
Thus by (8)
\begin{equation}
\Phi_1^n\Phi_n^{-1}=\prod_t \varphi_{t'}^{mE(\frac{nt}{m})}=\varepsilon\alpha^m.
\tag{10}
\end{equation}
Accordingly the quotient
\[
\left(\frac{\prod_t \varphi_{t'}^{E(\frac{nt}{m})}}{\alpha}\right)^m
\]
is a unit functional, and since it is at the same time the $m$th power of a functional in $\Omega_m$, the base of this power is also a unit, that is,
\[
\prod_t \varphi_{t'}^{E(\frac{nt}{m})}=\varepsilon\alpha,
\]
and this product belongs to the principal class. But by our assumption
\[
\varphi_{t'}\ \text{is equivalent to}\ \varphi^{t'},
\]
and hence
\[
\prod_t \varphi_{t'}^{E(\frac{nt}{m})}\ \text{is equivalent to}\ \varphi^{\sum t'E(\frac{nt}{m})},
\]
and this holds for every arbitrary $n$ not divisible by $q$.

If therefore we can prove that $n$ can be so chosen that also the sum
\[
S_n=\sum_t t'E\left(\frac{nt}{m}\right)
\]
is not divisible by $q$, then theorem a) is proved.

We first want to make a transformation of the sum $S_n$, by which the question is led back to the much simpler case in which $m$ is a prime number.

If $m$ is a higher than first power of $q$, we set
\[
m=qm',\qquad t=t_1+qt_2,
\qquad
\begin{array}{l}
t_1=1,2,\ldots,q-1,\\
t_2=0,1,\ldots,m'-1.
\end{array}
\]
Here $m'$ is still a power of $q$; the summation letter $t_1$ runs through the number sequence $1,2,\ldots,q-1$, and $t_2$ through the number sequence $0,1,2,\ldots,m'-1$. We determine $t_1'$ from the congruence
\[
t_1t_1'\equiv1\pmod q
\]
and obtain
\[
t'\equiv t_1'\pmod q.
\]
Then
\begin{equation}
S_n=\sum_t t'E\left(\frac{tn}{m}\right)
\equiv
\sum_{t_1}t_1'\sum_{t_2}E\left(\frac{nt_2}{m'}+\frac{nt_1}{m}\right)
\pmod q.
\tag{11}
\end{equation}
We now assume $n<q$, hence also $nt_1<m$, so that $nt_1:m$ is a proper fraction. Then
\[
\begin{array}{ll}
\text{a)}& E\left(\dfrac{nt_2}{m'}+\dfrac{nt_1}{m}\right)=E\left(\dfrac{nt_2}{m'}\right),\\[1.2em]
\text{b)}& \text{or}=E\left(\dfrac{nt_2}{m'}\right)+1,
\end{array}
\]
and case b) occurs only if between
\[
\frac{nt_2}{m'}\quad\text{and}\quad \frac{nt_2}{m'}+\frac{nt_1}{m}
\]
there lies an integer. This is the case only if the difference between $nt_2:m'$ and the next greater integer is smaller than $nt_1:m$, that is, if
\begin{equation}
E\left(\frac{nt_2}{m'}\right)+1-\frac{nt_2}{m'}<\frac{nt_1}{m}.
\tag{12}
\end{equation}
If in the expression on the left side of (12)
\begin{equation}
E\left(\frac{nt_2}{m'}\right)+1-\frac{nt_2}{m'}
\tag{13}
\end{equation}
we let $t_2$ run through all its values, we obtain positive rational fractions not exceeding unity, with denominator $m'$, no two of which are equal to one another. For if two of the values arising from (13) by substitution of $t_2',t_2''$ for $t_2$ were equal, then $n(t_2'-t_2'')/m'$ would have to be an integer, which, since $n$ is relatively prime to $m'$ and $t_2'-t_2''$ is smaller than $m'$, cannot be. Thus the expressions (13) run in some order through the numerical values
\begin{equation}
\frac1{m'},\ \frac2{m'},\ldots,\frac{m'-1}{m'},\ 1,
\tag{14}
\end{equation}
and if $a$ denotes an arbitrary positive fraction, then $E(m'a)$ is the number of numbers of the sequence (14) which are not greater than $a$. By (12), therefore, the number of values of $t_2$ for which case b) occurs is equal to $E(nt_1/q)$, and there results
\[
\sum_{t_2}E\left(\frac{nt_2}{m'}+\frac{nt_1}{m}\right)
=
\sum_{t_2}E\left(\frac{nt_2}{m'}\right)+E\left(\frac{nt_1}{q}\right).
\]
To form $S_n$, we multiply this equation by $t_1'$ and sum. This gives
\[
S_n\equiv \sum_{t_1}t_1'\sum_{t_2}E\left(\frac{nt_2}{m'}\right)+\sum_{t_1}t_1'E\left(\frac{nt_1}{q}\right)\pmod q.
\]
Here now $\sum t_1'=\sum t_1\equiv0\pmod q$, because the $t_1$ pairwise give the sum $q$, and it follows that
\begin{equation}
S_n\equiv\sum_{t_1}t_1'E\left(\frac{nt_1}{q}\right)\pmod q.
\tag{15}
\end{equation}
The sum standing here on the right side is formed in the same way as $S_n$ itself in the case where $m$ is a prime number. For this case, however, there results
\[
\frac{nt_1}{q}=E\left(\frac{nt_1}{q}\right)+\left(\frac{nt_1}{q}\right),
\]
\[
\frac{(q-n)t_1}{q}=E\left(\frac{(q-n)t_1}{q}\right)+\left(\frac{(q-n)t_1}{q}\right),
\]
and from this, by addition,
\[
t_1=E\left(\frac{nt_1}{q}\right)+E\left(\frac{(q-n)t_1}{q}\right)+\left(\frac{nt_1}{q}\right)+\left(\frac{(q-n)t_1}{q}\right),
\]
and here the sum of the two positive proper fractions, which nevertheless must be an integer, can have only the value 1. It follows therefore that
\[
E\left(\frac{nt_1}{q}\right)+E\left(\frac{(q-n)t_1}{q}\right)=t_1-1,
\]
and from this, by multiplication with $t_1'$ and summation over all values of $t_1$ from 1 to $q-1$,
\[
\sum_{t_1}t_1'E\left(\frac{nt_1}{q}\right)+\sum_{t_1}t_1'E\left(\frac{(q-n)t_1}{q}\right)\equiv -1\pmod q.
\]
Consequently both sums on the left side certainly cannot be divisible by $q$. Then expression (15) shows that of the two sums $S_n$ and $S_{q-n}$ one certainly is indivisible by $q$.

This, however, was to be proved.

\srcpage{666}
\section*{\S\,183. Proof of the Second Auxiliary Lemma for an Odd $m$.}

The second lemma, which by \S\,181 remains to be proved, is the following:

\begin{enumerate}[label=\arabic*.,leftmargin=2.2em,start=2]
\item If $\varepsilon_n$ is a system of conjugate numerical units of the field $\Omega_m$ which satisfies the condition that
\begin{equation}
\varepsilon_n\varepsilon_1^{-n}=\mathscr E^m
\tag{1}
\end{equation}
is the $m$th power of a unit in $\Omega_m$, then the $\varepsilon_n$ themselves are $m$th powers of units in $\Omega_m$, multiplied by a root of unity of the form $r^{nk}$.
\end{enumerate}

This proof too is quite elementary for the case of an odd $m$. We therefore first consider this case here.

By \S\,178, 1., we can represent every unit of the field $\Omega_m$ in the form
\begin{equation}
\varepsilon=\pm r^k\mathscr E(r),\qquad
\varepsilon_n=\pm r^{nk}\mathscr E(r^n),
\tag{2}
\end{equation}
where $\mathscr E(r)$ is a real unit of the field $\Omega_m$. Thus $\mathscr E(r)$ satisfies the condition
\begin{equation}
\mathscr E(r)=\mathscr E(r^{-1}),
\tag{3}
\end{equation}
from which, by (2), follows:
\begin{equation}
\varepsilon_1\varepsilon_{-1}=\mathscr E(r)^2;
\tag{4}
\end{equation}
but from (1), putting $n=-1$, one finds
\[
\varepsilon_1\varepsilon_{-1}=\mathscr E^m,
\]
that is, the left side of (4) is the $m$th power of a unit in $\Omega_m$, and hence also
\begin{equation}
\mathscr E(r)^2=\mathscr E^m
\tag{5}
\end{equation}
is the $m$th power of such a unit.

Since now $m$ is odd, the rational integer $x$ can be so determined that
\begin{equation}
2x+m=1,
\tag{6}
\end{equation}
and from this follows
\[
\mathscr E(r)=\mathscr E(r)^{2x}\mathscr E(r)^m.
\]
Thus if, according to (5), we set
\[
\mathscr E^x\mathscr E(r)=e(r),
\]
then follows
\begin{equation}
\mathscr E(r)=e(r)^m.
\tag{7}
\end{equation}
By (7), however, Theorem 2 is proved, and at the same time, for an odd $m$, the whole Theorem I.

This conclusion too fails for an even $m$, because then equation (6) is no longer soluble.

\srcpage{667}
\section*{\S\,184. Preliminary Remarks on the Case of an Even $m$.}

For the case where $m$ is a power of 2, we take an entirely different path, about which some preliminary remarks may first find place here.

We have already seen in \S\,182 that the first lemma is proved if we can prove that some power of $\varphi$, whose exponent is not divisible by $q$, belongs to the principal class. In doing this, of the two properties of the functional $\varphi$ only the first is used, namely that the $m$th power of $\varphi$ is contained in the principal class.

Now by the general theorems in \S\,154 we know in every field an exponent $h$, namely the number of ideal classes, for which $\varphi^h$ belongs to the principal class whenever $\varphi$ is any functional in this field.

Thus if we can prove that the class number $h$ in the field $\Omega_m$, when $m$ is a power of 2, is an odd number, then by this alone, without anything further, Lemma 1 is proved.

This is to be the aim of our considerations in the next sections.

With respect to the second lemma the matter is similar. Here, from the assumptions in \S\,183, 2., one can derive, exactly as above, formula \S\,183, (5), from which, however, it can only be concluded that $\mathscr E(r)$ is the $\frac12m$th power of a real unit $e(r)$.

Thus, assuming the first lemma as proved, formula \S\,181, (14) gives:
\begin{equation}
(r^n,x_0)=\rho_n\sqrt{e(r^n)}\,\alpha_n\prod_p (r^{m_2n},\eta),
\tag{1}
\end{equation}
where $\rho_n$ is a root of unity of degree $m^2$, and $\alpha_n$ is a number in $\Omega_m$, and it would still have to be proved that $e(r^n)$ is the square of a unit in $\Omega_m$.

Since $e(r)$ is a real unit, we have
\begin{equation}
e(r^n)=e(r^{-n}),
\tag{2}
\end{equation}
and we shall further stipulate that the sign of the root is to be so determined (which, by a suitable assumption about $\rho_n$, we can always assume) that
\begin{equation}
\sqrt{e(r^n)}=\sqrt{e(r^{-n})}.
\tag{3}
\end{equation}
By the assumption, the product $(r^n,x_0)(r^{-n},x_0)$ is a number of the field $\Omega_m$ which is not changed by the substitution $(r,r^{-1})$, that is, a real number. Further, by \S\,174, (5),
\[
(r^{m_2n},\eta)(r^{-m_2n},\eta)=\pm p,
\]
so it is likewise real. Similarly $\alpha_n\alpha_{-n}$ is real, and from (1) it follows that
\[
\rho_n\rho_{-n}=\frac{(r^n,x_0)(r^{-n},x_0)}{e(r^n)\alpha_n\alpha_{-n}\prod_p (r^{m_2n},\eta)(r^{-m_2n},\eta)}
\]
is a real number, which, since it is a root of unity, can only be $=\pm1$.

Now in the equation resulting from this for $n=1$,
\[
(r,x_0)(r^{-1},x_0)=\pm e(r)\alpha_1\alpha_{-1}\prod_p(\pm p),
\]
we can carry out every substitution $(r,r^n)$, from which it follows that the sign of $\rho_n\rho_{-n}$ is independent of $n$, hence that
\begin{equation}
\rho_n\rho_{-n}=\rho_1\rho_{-1}=\pm1.
\tag{4}
\end{equation}
Now we further derive from (1):
\begin{equation}
\rho_n\rho_1^{-1}\sqrt{e(r^n)}\sqrt{e(r)}^{-n}=\frac{(r^n,x_0)(r,x_0)^{-n}}{\alpha_n\alpha_1^{-1}\prod_p(r^{m_2n},\eta)(r^{m_2},\eta)^{-n}},
\tag{5}
\end{equation}
and the right side is a number of the field $\Omega_m$. But the left side shows that it is a unit, and therefore, if $\mathscr E(r)$ denotes a real unit, we can set
\begin{equation}
\rho_n\rho_1^{-1}\sqrt{e(r^n)}\sqrt{e(r)}^{-n}=r^k\mathscr E(r).
\tag{6}
\end{equation}
If we substitute this value into formula (5) and make the substitution $(r,r^{-1})$, there results
\begin{equation}
\rho_{-n}\rho_{-1}^{-1}\sqrt{e(r^n)}\sqrt{e(r)}^{-n}=r^{-k}\mathscr E(r),
\tag{7}
\end{equation}
and by multiplication of (6) and (7), with regard to (4),
\begin{equation}
e(r^n)e(r)^{-n}=\mathscr E(r)^2.
\tag{8}
\end{equation}

In our considerations on the class number the following theorem will still result:
\begin{center}
A unit $e(r)$ of the real field $H_m$ which is positive with all its conjugates is the square of a unit in the field $H_m$.
\end{center}
Formula (8), however, shows that the unit $\pm e(r)$ has this property, and hence that $e(r)$ (since also $-1=i^2$ is the square of a unit) is the square of a unit of the field $\Omega_m$.

Thus the question has been reduced to the proof of the two mentioned theorems, which arises as the capstone of a long chain of considerations, but one extremely rich in interesting relations, to which the next sections are to be devoted.

\srcpage{670}
\section*{Twenty-third Section. Class Number.}

\section*{\S~185. Dirichlet's theorem on the units.}

The investigations on the number of ideal classes in an algebraic field, which are now to give us the necessary supplements to the proofs of the preceding section, make use of the methods that Dirichlet introduced into number theory.\footnote{Dirichlet, \textit{Recherches sur diverses applications de l'analyse infinitésimale à la théorie des nombres}. Crelle's Journal, vols.~19, 21. Dirichlet's Works, vol.~I (1839, 1840). (The posthumous papers of Gauss, ``De nexu inter multitudinem classium etc.'' Gauss' Works, vol.~II, also belong here.) The application of these principles to cyclotomic numbers was first made by Kummer [Crelle's Journal, vols.~35, 40 (1847, 1850); Liouville's Journal, vol.~16 (1851)]. The general theory of units was likewise founded by Dirichlet [Dirichlet's Works, vol.~I, pp.~622, 633, 639 (1840, 1842, 1846)]. The general theory of units and the determination of the class numbers was given by Dedekind: Dirichlet-Dedekind, \textit{Vorlesungen über Zahlentheorie}, 4th ed., \S~183 ff. Minkowski, \textit{Geometrie der Zahlen}.}

They are no longer purely algebraic, insofar as transcendental functions and numbers, such as the natural logarithm and the number $\pi$, also occur in them. Use is also made of passages to the limit, as in integral calculus, from which we have already drawn a fine result in \S~156.

The investigations with which we first have to deal are quite general, that is, they hold for every algebraic number field, and it offers no advantage for simplicity to restrict them to special fields, such as the cyclotomic fields.

\srcpage{671}
Thus let $\Omega$ be a field of the $n$th degree, and let
\begin{equation}
\Omega_1,\Omega_2,\ldots,\Omega_n
\tag{1}
\end{equation}
be the conjugate fields. The irreducible equation of the $n$th degree whose roots determine these conjugate fields for us will in general have both real and imaginary roots; the latter, however, always occur pairwise, so that to every imaginary field $\Omega^{(i)}$ there belongs another whose numbers are conjugate imaginary to those of $\Omega^{(i)}$, and which therefore arises from $\Omega^{(i)}$ by the substitution $(i,-i)$.

We combine such imaginary pairs into one unit and denote by $\nu$ the number of real fields and of pairs of imaginary fields in the sequence (1). If all conjugate fields are real, then $n=\nu$; if they are all imaginary, then $n=2\nu$.

Except in the case of the rational field, $\nu$ is equal to $1$ only in the case of an imaginary quadratic field. Otherwise $\nu$ is always greater than $1$.

The number of imaginary pairs is $n-\nu$, and the number of real fields is $2\nu-n$.

In the determinant by which, in \S~145, the square root of the ground number, $\sqrt{\Delta}$, is defined, the substitution $(i,-i)$ interchanges two conjugate imaginary rows at a time, and $\sqrt{\Delta}$ therefore changes its sign precisely when this number is odd, and otherwise not. In the first case $\sqrt{\Delta}$ is imaginary, in the second real. Hence it follows that the ground number $\Delta$ of the field $\Omega$ is positive or negative according as $n-\nu$ is even or odd.

If from two conjugate imaginary fields we retain only one, there results the sequence of conjugate fields
\[
\Omega_1,\Omega_2,\ldots,\Omega_\nu,
\]
to which we assign a system of signs
\[
\delta_1,\delta_2,\ldots,\delta_\nu
\]
with the convention that $\delta_s=1$ if $\Omega_s$ is real, and $=2$ if $\Omega_s$ is imaginary, so that $\sum \delta_s=n$.

Let $\eta$ denote a non-zero number of the field $\Omega$ and let
\[
\eta_1,\eta_2,\ldots,\eta_n
\]
be the conjugate numbers. To this number system we assign another number system
\[
\lambda_1,\lambda_2,\ldots,\lambda_\nu,
\]
\srcpage{672}
which we define by the equation
\begin{equation}
\lambda_s=\delta_s\log |\eta_s|,
\tag{2}
\end{equation}
where $|\eta_s|$ denotes the absolute value of $\eta_s$, and $\log|\eta_s|$ denotes the real natural logarithm of the positive number $|\eta_s|$. If $\Omega_s$ is real and the sign $\pm$ is chosen so that $\pm\eta_s$ is positive, then
\[
\lambda_s=\log(\pm\eta_s).
\]
But if $\eta_s,\eta_s'$ is an imaginary pair, then
\[
\lambda_s=\log\eta_s\eta_s',
\]
and to $\eta_s$ and $\eta_s'$ belongs the same $\lambda_s$, so that the number of different $\lambda_s$ is equal to $\nu$. From this determination one further obtains
\begin{equation}
\lambda_1+\lambda_2+\cdots+\lambda_\nu=\log N_a(\eta),
\tag{3}
\end{equation}
where, as before, $N_a$ denotes the absolute norm.

We shall call the numbers $\lambda_1,\lambda_2,\ldots,\lambda_\nu$ the conjugate logarithms of the number $\eta$.

If $\eta$ is an integer of the field $\Omega$, then its absolute norm is a natural number, which is equal to $1$ only when $\eta$ is a unit; and from (3) it follows that:

\textbf{1. The sum of the conjugate logarithms of an integer $\eta$ is positive, and is equal to zero only when $\eta$ is a unit.}

Let $\omega_1,\omega_2,\ldots,\omega_n$ denote a basis of the integers in $\Omega$ (a minimal basis of $\Omega$, \S~145), and let $\omega_{1,s},\omega_{2,s},\ldots,\omega_{n,s}$, for $s=1,2,\ldots,n$, denote the conjugate numbers. Then all integers of the field $\Omega_s$ can be represented in the form
\begin{equation}
\eta_s=x_1\omega_{1,s}+x_2\omega_{2,s}+\cdots+x_n\omega_{n,s},
\tag{4}
\end{equation}
with rational integral coefficients $x_1,x_2,\ldots,x_n$. The determinant of this system,
\[
\sum \pm \omega_{1,1}\omega_{2,2}\cdots\omega_{n,n}=\sqrt{\Delta},
\]
is the square root of the discriminant $\Delta$ of the field, and hence is always different from zero. Thus the system (4) can be solved for the unknowns $x_1,x_2,\ldots,x_n$, and if we now require that the absolute values of the numbers $\eta_s$ are not to exceed a certain bound, then these solutions give bounds for the rational integers $x_i$.

\srcpage{673}
This simple remark gives us the following important theorem:

\textbf{2. In an algebraic field $\Omega$ there is only a finite number of integers $\eta$ of the property that the absolute values of the conjugate numbers $\eta_s$ lie below a given bound.}

We now make use of the general theorem \S~155, (25), according to which, if $f(x_1,x_2,\ldots,x_n)$ is any positive quadratic form in $n$ variables with determinant $D$, the variables $x_i$ can be determined as rational integers, not all zero, so that
\[
f(x_1,x_2,\ldots,x_n)< n\sqrt[n]{0.404D},
\]
that is, becomes smaller than a number depending only on the determinant $D$.

We apply this theorem, as in \S~156, to the form
\[
f(x_1,x_2,\ldots,x_n)=\frac{|\eta_1|^2}{c_1^2}+\frac{|\eta_2|^2}{c_2^2}+\cdots+\frac{|\eta_n|^2}{c_n^2},
\]
where the $\eta_s$ are the linear forms (4), and in the case of a real $\eta_s$
\[
|\eta_s|^2=(\eta_s)^2,
\]
in the case of an imaginary pair $\eta_s,\eta_s'$
\[
|\eta_s|^2=|\eta_s'|^2=\eta_s\eta_s'.
\]
In the latter case $\eta_s\eta_s'$ is the sum of two real squares, and hence, if the $c_s$ are taken to be real, $f$ is a positive form. We further stipulate, for the hitherto arbitrary real quantities $c_s$, that if $\eta_s,\eta_s'$ form a conjugate imaginary pair, then
\begin{equation}
c_s=c_s'
\tag{5}
\end{equation}
is to hold; besides this we subject the $c_s$ to the condition
\begin{equation}
c_1c_2\cdots c_n=1.
\tag{6}
\end{equation}
We form the determinant of the function $f$, as in \S~156, by first regarding it as a form in the variables $\eta_1,\eta_2,\ldots,\eta_n$, and the determinant of this form is, because of (6), equal to $\pm1$. The determinant of $f$ with respect to the variables $x_i$ is therefore (Vol.~I, \S~56), apart from the sign, equal to the
\srcpage{674}
square of the determinant of the linear forms $\eta_s$, that is, to the ground number of the field $\Omega$.

It follows that one can find a number $A$, independent of the $c_s$ and determined only by the field $\Omega$, of such a kind that, whatever the $c_s$ may be, $f(x_1,x_2,\ldots,x_n)$ for integral non-zero $x$ drops below $A$. Since $f$ is a sum of positive summands, the same holds for each of these summands, and the following theorem is thereby proved:

\textbf{3. If $c_s$ is any system of real positive numbers satisfying the conditions (5), (6), then there is a real number $A$, determined by the field $\Omega$ alone, such that one can always determine an integer $\eta$ of the field $\Omega$ which, together with its conjugate numbers $\eta_s$, satisfies the conditions}
\begin{equation}
|\eta_s|<c_sA.
\tag{7}
\end{equation}

Incidentally it follows from this that in every algebraic field $\Omega$, except the rational and the imaginary quadratic field, one can find non-zero integers whose absolute value falls below any given bound.

We shall transform the expression of theorem 3 somewhat by introducing the conjugate logarithms of $\eta$. Put $\log A=k$, so that $k$ is a real number determined by $\Omega$, and further
\begin{equation}
c_s=e^{\frac{\gamma_su}{\delta_s}},
\tag{8}
\end{equation}
where $u$ may be any real quantity, and where the numbers $\gamma_s$ satisfy, because of (5) and (6), the condition
\begin{equation}
\gamma_1+\gamma_2+\cdots+\gamma_\nu=\sum_{s=1}^{\nu}\gamma_s=0.
\tag{9}
\end{equation}
Since, moreover,
\[
|\eta_s|=e^{\frac{\lambda_s}{\delta_s}},
\]
it follows from (7) and (8) that
\begin{equation}
\lambda_s-\gamma_su<k\delta_s.
\tag{10}
\end{equation}
The sum of the $\nu$ expressions on the left side of this inequality is, by (9), equal to $\sum\lambda_s$, hence, by theorem 1, is not negative; and we find
\[
0\leq \sum_s(\lambda_s-\gamma_su)<kn.
\]

\srcpage{675}
It follows from this, in view of (10), that a term of this sum, $\lambda_s-\gamma_su$, cannot be smaller than $-(n-\delta_s)k$, for otherwise the total sum would be negative; and if, for simplicity, we take the lower bound somewhat smaller,
\begin{equation}
-nk\delta_s<\lambda_s-\gamma_su<k\delta_s.
\tag{11}
\end{equation}

It is thereby proved that for every given system of numbers $u,\gamma_s$, provided only that condition (9) is fulfilled, there is an integer $\eta$ of the field $\Omega$ which, with its conjugate numbers, satisfies the limiting conditions (11).

Now let $g_s$ be a system of real quantities of which we demand only that
\[
\gamma_1g_1+\gamma_2g_2+\cdots+\gamma_\nu g_\nu=\alpha
\]
be different from zero. By (9) this excludes the case where all the $g_s$ are equal to one another; but if they are not, then for every given system of the $g_s$ there will be infinitely many systems $\gamma_s$ that satisfy the requirement (9).

If now $k\sum\delta_sg_s=g$ is put, then from (11), by multiplication with $g_s$ and summation, one obtains
\begin{equation}
-ng+\alpha u<\sum g_s\lambda_s<g+\alpha u.
\tag{12}
\end{equation}

According to this limiting condition, keeping $g_s$ and $\alpha$ fixed, we now determine a sequence of integers
\begin{equation}
\eta,\eta',\eta'',\eta''',\ldots,
\tag{13}
\end{equation}
by making, for $u$, a sequence of assumptions
\[
u,u',u'',u''',\ldots,
\]
which are more precisely determined as follows.

First choose $u$ so that
\[
-ng+\alpha u=a,
\qquad
 g+\alpha u=a'
\]
are positive; then put
\[
\delta=\frac{a'-a}{\alpha}=\frac{(n+1)g}{\alpha},
\]
and put
\[
u'=u+\delta,
\quad
u''=u'+\delta,
\quad
u'''=u''+\delta,
\quad\ldots,
\]
\[
a+\alpha\delta=a',
\quad a'+\alpha\delta=a'',
\quad a''+\alpha\delta=a''',
\quad\ldots,
\]
\srcpage{676}
and from (12) we thus obtain, if $\lambda_s',\lambda_s'',\ldots$ are the conjugate logarithms of $\eta',\eta'',\ldots$,
\begin{equation}
\begin{array}{rcl}
a&<&\sum g_s\lambda_s<a',\\
a'&<&\sum g_s\lambda_s'<a'',\\
a''&<&\sum g_s\lambda_s''<a''',\\
&&\cdots\cdots,
\end{array}
\tag{14}
\end{equation}
and the sequence of these inequalities can be continued without bound.

Moreover all these numbers $\eta,\eta',\eta'',\ldots$ have, by (6) and (7), the property that their absolute norms $N,N',N'',\ldots$ remain below a fixed finite bound $e^{nk}$.

The number of possible values of $N,N',N'',\ldots$ is therefore finite, certainly not greater than the greatest integer contained in $e^{nk}$. Likewise the number of all possible residues of the numbers $x_1,x_2,\ldots,x_n$ [in (4)] modulo all the moduli $N,N',N'',\ldots$ is finite; and from this it follows that, if we continue the sequence of numbers (13) far enough, two distinct numbers $\eta^{(h)},\eta^{(k)}$ must occur in the sequence for which $N^{(h)}=N^{(k)}$ and the residues of $x_1,x_2,\ldots,x_n$ modulo $N^{(h)}$ are exactly the same.

But from (14), if $h>k$ is assumed, it follows that
\begin{equation}
\sum g_s\bigl(\lambda_s^{(h)}-\lambda_s^{(k)}\bigr)>0.
\tag{15}
\end{equation}
If $N$ is the absolute norm of these two numbers $\eta^{(h)},\eta^{(k)}$, then, because of the agreement of the residues of the $x$, the difference $\eta^{(h)}-\eta^{(k)}$ is divisible by $N$. On the other hand, by \S~138 (13), the number $N$ is divisible by $\eta^{(h)}$, and consequently $\eta^{(h)}-\eta^{(k)}$ is also divisible by $\eta^{(h)}$. It follows that $\eta^{(k)}$ is also divisible by $\eta^{(h)}$, and likewise $\eta^{(h)}$ by $\eta^{(k)}$. The two numbers are therefore associated, and if we put
\[
\frac{\eta^{(h)}}{\eta^{(k)}}=\varepsilon,
\]
then $\varepsilon$ is a unit of the field $\Omega$.

Moreover, if, denoting by $|\varepsilon_s|$ the absolute value of $\varepsilon_s$, we put
\begin{equation}
\delta_s\log|\varepsilon_s|=l_s,
\tag{16}
\end{equation}
then from (15) we obtain
\begin{equation}
g_1l_1+g_2l_2+\cdots+g_\nu l_\nu>0.
\tag{17}
\end{equation}

\srcpage{677}
This proves the theorem:

\textbf{4. If $g_1,g_2,\ldots,g_\nu$ is any system of real numbers which are not all equal to one another, then in every algebraic field $\Omega$ one can determine a unit $\varepsilon$ such that the sum}
\[
g_1l_1+g_2l_2+\cdots+g_\nu l_\nu
\]
\textbf{is different from zero, where $l_s$ is the system of conjugate logarithms of $\varepsilon$.}

This important theorem, which is the foundation for what follows, is due, though in a somewhat changed form, to Dirichlet. This formulation was given by Minkowski.

\section*{\S~186. Systems of independent units and exponent systems of the units.}

If in the theorem just proved we put $g_1=1$, $g_2=0,\ldots,g_\nu=0$, which is allowed if from now on we assume $\nu>1$, it follows that in $\Omega$ there is a unit $\varepsilon$ for which one of the conjugate logarithms, $l_1$, is different from zero. Since every power of $\varepsilon$ has the same property, there are also infinitely many such units.

This theorem now permits a very important generalization:

\textbf{5. If $s\leq \nu-1$, then in the field $\Omega$ one can determine a system of units with the associated conjugate logarithms}
\begin{equation}
\begin{array}{cccc}
\varepsilon_1: & l_{1,1},&l_{1,2},&\ldots,l_{1,\nu},\\
\varepsilon_2: & l_{2,1},&l_{2,2},&\ldots,l_{2,\nu},\\
&\multicolumn{3}{c}{\cdots\cdots\cdots}\,\\
\varepsilon_s: & l_{s,1},&l_{s,2},&\ldots,l_{s,\nu},
\end{array}
\tag{1}
\end{equation}
\textbf{in such a way that any one of the $s$-rowed determinants of the matrix of the $l_{h,k}$, for instance}
\[
L_s=\sum \pm l_{1,1}l_{2,2}\cdots l_{s,s},
\]
\textbf{is different from zero.}

For $s=1$ the stated theorem is correct by the remark made at the beginning. We therefore assume that it
\srcpage{678}
has been proved for $s-1$ and derive it generally by complete induction. For this purpose we arrange the determinant $L_s$ according to the elements of the last row, and write it thus:
\begin{equation}
L_s=g_1l_{s,1}+g_2l_{s,2}+\cdots+g_sl_{s,s},
\tag{2}
\end{equation}
where $g_s=L_{s-1}$ and hence, by the assumption made, can be taken to be different from zero. If we substitute these values of $g$ in the formula of theorem 4, \S~185, taking $g_{s+1}=0,\ldots,g_\nu=0$, while $g_s$ is not equal to zero, then, as long as $s<\nu$, certainly not all $g_1,g_2,\ldots,g_\nu$ are equal to one another; and it follows that $\varepsilon_s$ can be so chosen that $L_s$ is different from zero, as was to be proved.

This mode of inference cannot be extended to $s=\nu$, and the determinant $L_\nu$ must indeed always be zero, since for every unit $\sum_{1}^{\nu} l_s$ vanishes.

For the case $s=\nu-1$ we shall now denote the determinant $L_s$ by
\begin{equation}
L_{\nu-1}=L(\varepsilon_1,\varepsilon_2,\ldots,\varepsilon_{\nu-1}),
\tag{3}
\end{equation}
and call a system of units $\varepsilon_1,\varepsilon_2,\ldots,\varepsilon_{\nu-1}$ for which this determinant is different from zero a system of independent units.

The absolute value $L$ of the determinant $L(\varepsilon_1,\varepsilon_2,\ldots,\varepsilon_{\nu-1})$ is called the regulator of the system $\varepsilon_1,\varepsilon_2,\ldots,\varepsilon_{\nu-1}$.\footnote{Dirichlet-Dedekind, \textit{Vorlesungen über Zahlentheorie}, 4th ed., \S~183.}

If in the determinant
\begin{equation}
\left|\begin{array}{cccc}
l_{1,1},&l_{1,2},&\ldots,&l_{1,\nu}\\
\cdots&\cdots&\cdots&\cdots\\
l_{\nu-1,1},&l_{\nu-1,2},&\ldots,&l_{\nu-1,\nu}\\
x_1,&x_2,&\ldots,&x_\nu
\end{array}\right|,
\tag{4}
\end{equation}
in which $x_1,x_2,\ldots,x_\nu$ are arbitrary quantities, we add all columns to the last one, then, with regard to the relations $\sum_{1}^{\nu}l_{s,i}=0$, we obtain its value equal to
\[
\pm(x_1+x_2+\cdots+x_\nu)L.
\]
If, therefore, we take all the $x$ except for an arbitrary one
\srcpage{679}
to be $=0$, and the remaining one to be $=1$, then from (3) there results one of the $(\nu-1)$-rowed determinants of the matrix
\[
\begin{array}{ccc}
l_{1,1},&\ldots,&l_{1,\nu}\\
\cdots&\cdots&\cdots\\
l_{\nu-1,1},&\ldots,&l_{\nu-1,\nu}
\end{array}
\]
and all these determinants therefore have, apart from the sign, the same value. It is therefore immaterial which one of the $\nu$ conjugate values is omitted in forming the regulator $L$.

If $\varepsilon_1,\varepsilon_2,\ldots,\varepsilon_{\nu-1}$ is an independent system of units, and if $\lambda_1,\lambda_2,\ldots,\lambda_\nu$ are the conjugate logarithms of any unit $\varepsilon$ in $\Omega$, then the numbers $\xi_1,\xi_2,\ldots,\xi_{\nu-1}$ can always and uniquely be determined so that
\begin{equation}
\begin{array}{rcl}
\xi_1l_{1,1}+\xi_2l_{2,1}+\cdots+\xi_{\nu-1}l_{\nu-1,1}&=&\lambda_1,\\
\xi_1l_{1,2}+\xi_2l_{2,2}+\cdots+\xi_{\nu-1}l_{\nu-1,2}&=&\lambda_2,\\
\multicolumn{3}{c}{\cdots\cdots\cdots}\,\\
\xi_1l_{1,\nu}+\xi_2l_{2,\nu}+\cdots+\xi_{\nu-1}l_{\nu-1,\nu}&=&\lambda_\nu
\end{array}
\tag{5}
\end{equation}
holds. For the first $\nu-1$ of these equations form a system of linear equations with non-vanishing determinant, and the $\nu$th equation follows from the others because the sum of the conjugate logarithms of every unit vanishes.

We call the system of numbers $\xi_1,\xi_2,\ldots,\xi_{\nu-1}$ the exponent system of the unit $\varepsilon$ with respect to the system $\varepsilon_1,\varepsilon_2,\ldots,\varepsilon_{\nu-1}$, and it remains to prove:

\textbf{6. That the exponents of every unit $\varepsilon$ are rational numbers whose denominators do not exceed a certain finite bound, and which therefore can all be assumed to have a common denominator depending only on the system $\varepsilon_i$.}

To see this, one must consider the following:

1) If we regard the quantities $\xi_1,\xi_2,\ldots,\xi_{\nu-1}$ as variables and let each of them, independently of the others, run from 0 to 1, then the left sides of (5) remain enclosed within finite bounds depending only on the units $\varepsilon_1,\varepsilon_2,\ldots,\varepsilon_{\nu-1}$. Within these bounds, therefore, the conjugate logarithms $\lambda$ of the unit $\varepsilon$ must also remain, as long as
\srcpage{680}
the exponents $\xi_i$ remain restricted to non-negative proper fractional values; and hence, from the definition of the conjugate logarithms, one obtains an upper bound for $\varepsilon$ and its conjugates.

But by theorem 2, \S~185, there are in $\Omega$ only finitely many integers, and therefore still more only finitely many units, which in absolute value, with all their conjugates, lie below a finite bound.

If we now call a unit whose exponents lie between the bounds 0 and 1, including the lower and excluding the upper bound, a unit reduced with respect to the system $\varepsilon_1,\varepsilon_2,\ldots,\varepsilon_{\nu-1}$, then we have the theorem:

\textbf{There are only finitely many units in $\Omega$ which are reduced with respect to an independent system $\varepsilon_1,\varepsilon_2,\ldots,\varepsilon_{\nu-1}$.}

2) The units reproduce themselves by multiplication and division. If
\[
\begin{array}{c}
\xi_1',\xi_2',\ldots,\xi_{\nu-1}',\\
\xi_1'',\xi_2'',\ldots,\xi_{\nu-1}''
\end{array}
\]
are the exponents of two units $\varepsilon',\varepsilon''$, then the sums and the differences
\[
\xi_1'\pm\xi_1'',\ \xi_2'\pm\xi_2'',\ldots,\xi_{\nu-1}'\pm\xi_{\nu-1}''
\]
are the exponents of the units $\varepsilon'\varepsilon''$ and $\varepsilon':\varepsilon''$.

This follows immediately from the theorem that the logarithm of a product or of a quotient is equal to the sum or the difference of the logarithms of the two components.

3) The units $\varepsilon_1,\varepsilon_2,\ldots$ themselves have the exponents
\[
1,0,\ldots,0;\quad 0,1,\ldots,0;\quad\ldots,
\]
and from 2) it follows that every system of integers, substituted for $\xi_1,\xi_2,\ldots,\xi_{\nu-1}$, gives the exponent system of a unit which can be formed by multiplication of powers of the $\varepsilon_1,\varepsilon_2,\ldots$.

4) From 2) and 3) it follows that an exponent system $\xi_1,\xi_2,\ldots,\xi_{\nu-1}$ of a unit remains an exponent system if all its elements are multiplied by one and the same integer, and that it also retains this character when its elements $\xi_i$ are increased or decreased by arbitrary integers. Thus, if we use the sign $(x)$ in the sense of \S~182, that it denotes the excess
\srcpage{681}
of the number $x$ over the greatest integer contained in $x$, the following is obtained:

If $\xi_1,\xi_2,\ldots,\xi_{\nu-1}$ is the exponent system of a unit, and $m$ is an integer, then also
\begin{equation}
(m\xi_1),(m\xi_2),\ldots,(m\xi_{\nu-1})
\tag{6}
\end{equation}
is the exponent system of a unit.

Now the quantities $(m\xi_i)$, when they are not zero, are positive proper fractions; and by 1), if the sequence of integers is continued sufficiently far, it must therefore occur that for two distinct values $m',m''$ of $m$ the number sequence (6) agrees. Since accordingly $m'\xi_i,m''\xi_i$ have the same excess over an integer, $(m'-m'')\xi_i$ is itself an integer; and so there is a rational integer $m$, at most equal to the number of exponent systems of reduced units, with the property that
\[
m\xi_1,m\xi_2,\ldots,m\xi_{\nu-1}
\]
become rational integers.

This number $m$ may change when the unit $\varepsilon$ is changed. But since all possible values of this number lie below a finite bound, there is a finite number in which all these values of $m$ are contained and which itself can be taken for $m$. Hence there is a certain positive integer $m$ which can be taken as denominator of all rational numbers which can occur as exponents of units. This proves theorem 6.

\section*{\S~187. Fundamental systems of units.}

We now choose, instead of the system of independent units $\varepsilon_i$, another system whose conjugate logarithms may be
\begin{equation}
\lambda_{i,1},\lambda_{i,2},\ldots,\lambda_{i,\nu}
\qquad i=1,2,\ldots,\nu-1.
\tag{1}
\end{equation}
By \S~186, (5),
\begin{equation}
\begin{array}{rl}
\lambda_{i,k}={}&\xi_{1,i}l_{1,k}+\xi_{2,i}l_{2,k}+\cdots+\xi_{\nu-1,i}l_{\nu-1,k},\\[2pt]
&i=1,2,\ldots,\nu-1;\quad k=1,2,\ldots,\nu,
\end{array}
\tag{2}
\end{equation}
where $\xi_{1,i},\xi_{2,i},\ldots,\xi_{\nu-1,i}$ denote the exponents of one of the new units with respect to the system $\varepsilon_1,\varepsilon_2,\ldots,\varepsilon_{\nu-1}$, which we may assume to be rational numbers with the common denominator $m$.

\srcpage{682}
But by the multiplication theorem for determinants,
\begin{equation}
\left|\begin{array}{ccc}
\lambda_{1,1},&\ldots,&\lambda_{\nu-1,1}\\
\cdots&\cdots&\cdots\\
\lambda_{1,\nu-1},&\ldots,&\lambda_{\nu-1,\nu-1}
\end{array}\right|
=
\left|\begin{array}{ccc}
\xi_{1,1},&\ldots,&\xi_{1,\nu-1}\\
\cdots&\cdots&\cdots\\
\xi_{\nu-1,1},&\ldots,&\xi_{\nu-1,\nu-1}
\end{array}\right|
\left|\begin{array}{ccc}
l_{1,1},&\ldots,&l_{\nu-1,1}\\
\cdots&\cdots&\cdots\\
l_{1,\nu-1},&\ldots,&l_{\nu-1,\nu-1}
\end{array}\right|,
\tag{3}
\end{equation}
or, if we put
\[
\Lambda_{\nu-1}=\sum\pm\lambda_{1,1}\lambda_{2,2}\cdots\lambda_{\nu-1,\nu-1}
\]
and observe that the determinant
\[
\sum\pm\xi_{1,1}\xi_{2,2}\cdots\xi_{\nu-1,\nu-1}
\]
is a rational number with denominator $m^{\nu-1}$,
\begin{equation}
\Lambda_{\nu-1}=\frac{a}{m^{\nu-1}}L(\varepsilon_1,\varepsilon_2,\ldots,\varepsilon_{\nu-1}),
\tag{4}
\end{equation}
where $a$ is a rational integer.

It follows that the newly introduced units form an independent system if and only if the determinant of the $\xi_{i,k}$, that is, the number $a$, is different from zero.

Now among a finite or infinite number of non-zero rational fractions with the same denominator there is always one of least absolute value, and consequently, by (4), we can choose the new system of independent units so that its regulator $\pm\Lambda_{\nu-1}$ becomes as small as possible. Such a system we call a fundamental system of units; and the minimum value of the regulator itself, that is, the regulator of a fundamental system, we call the regulator of the field. We now assume that our system $\varepsilon_1,\varepsilon_2,\ldots,\varepsilon_{\nu-1}$ itself is a fundamental system.

Under this assumption it can be proved that all exponents of units must be integers.

Indeed, if we assume that there exists a unit $\varepsilon$ whose exponents determined by the formulae \S~186, (5) are not all integers, then by \S~186, 4) there is also a unit $\varepsilon_0$ whose exponents are proper fractions and do not all vanish.

If by $\lambda_1,\lambda_2,\ldots,\lambda_{\nu-1}$ in the formulae (5), \S~186, we understand the system of conjugate logarithms of $\varepsilon_0$, then
\[
L(\varepsilon_0,\varepsilon_2,\ldots,\varepsilon_{\nu-1})=\xi_1L(\varepsilon_1,\varepsilon_2,\ldots,\varepsilon_{\nu-1}),
\]
\srcpage{683}
and if we therefore assume, which is plainly no restriction of generality, that $\xi_1$ is a non-zero proper fraction, then this formula contradicts the assumption that
\[
\varepsilon_1,\varepsilon_2,\ldots,\varepsilon_{\nu-1}
\]
is a fundamental system, because the determinant $L_{\nu-1}$ would be made smaller if $\varepsilon_1$ were replaced by $\varepsilon_0$.

Once we have a fundamental system, we can take any system of rational integers $\xi_1,\xi_2,\ldots,\xi_{\nu-1}$ as an exponent system and form a unit $\varepsilon$ with these exponents:
\[
\varepsilon=\varepsilon_1^{\xi_1}\varepsilon_2^{\xi_2}\cdots\varepsilon_{\nu-1}^{\xi_{\nu-1}}.
\]

It remains to answer the question how far a unit is determined by its exponent system.

Suppose that $\varepsilon',\varepsilon''$ are two units with the same exponent system. Then the exponent system of the unit $\varepsilon':\varepsilon''=\rho$ consists entirely of zeros, and consequently the conjugate logarithms of $\rho$ are all equal to zero. Therefore $\rho$ is an integer with the property that the absolute values of all numbers conjugate with $\rho$ are equal to $1$, and hence, as proved in \S~175, 2, it is a root of unity. Thus the following theorem is proved:

\textbf{1. There is in the field $\Omega$ a system of $\nu-1$ fundamental units}
\begin{equation}
\varepsilon_1,\varepsilon_2,\ldots,\varepsilon_{\nu-1},
\tag{5}
\end{equation}
\textbf{having the property that in the form}
\begin{equation}
\varepsilon=\rho\varepsilon_1^{\xi_1}\varepsilon_2^{\xi_2}\cdots\varepsilon_{\nu-1}^{\xi_{\nu-1}}
\tag{6}
\end{equation}
\textbf{all units of the field, each only once, are contained, when $\xi_1,\xi_2,\ldots,\xi_{\nu-1}$ run through all rational integers and $\rho$ runs through all roots of unity present in $\Omega$.}

It is now easy to derive all other fundamental systems from one fundamental system, by putting in (6), for the exponents, $\nu-1$ different systems of integers $\xi_{i,k}$ whose determinant is $=\pm1$. For if one puts
\[
\varepsilon_i'=\rho_i\varepsilon_1^{\xi_{1,i}}\varepsilon_2^{\xi_{2,i}}\cdots\varepsilon_{\nu-1}^{\xi_{\nu-1,i}},
\]
then from (3) it follows that
\[
L(\varepsilon_1',\varepsilon_2',\ldots,\varepsilon_{\nu-1}')=\pm L(\varepsilon_1,\varepsilon_2,\ldots,\varepsilon_{\nu-1}).
\]

\srcpage{684}
Thus both determinants have the minimum value, and both systems are fundamental systems.

The case $\nu=1$, which we excluded above, occurs, apart from the case of the rational field, in which only the two units $\pm1$ exist, in the case of the imaginary quadratic field. In this case, in the twentieth section, we represented the integers of the field in the form
\[
\frac{x+y\sqrt{-m}}{2},
\]
where $-m$ denotes a fundamental discriminant, so that $m\equiv0$ or $\equiv3\pmod 4$, and $x,y$ are both even or both odd. We obtain the units if we solve the equation
\[
x^2+my^2=4
\]
in all possible ways in rational integers. But this equation, if $m>4$, has only the two solutions $x=\pm2,y=0$, and thus in these cases, as in the field of rational numbers, there are only the two units $\pm1$. If $m=4$, one finds the four units $\pm1,\pm i$, and finally if $m=3$, the six units
\[
\pm1,
\quad \pm\frac{-1+i\sqrt3}{2},
\quad \pm\frac{-1-i\sqrt3}{2}.
\]

In the next simplest case, $n=2$, $\nu=2$, that is, in the real quadratic field, the theory of units coincides with the theory of Pell's equation treated in \S~127 of the first volume.
\end{document}
